\cotefile{29}
\scanpage{188}
Opérateur d'\uline{homotopie} pour $d''$

% The pencilled wave cancels the two upper lines. A large X crosses the draft below.
\[
\begin{aligned}
&\deleted{\overline W_n(G)=G_{n-1}\times G_{n-2}\times\cdots\times G_0}\\
&\deleted{\deleted{\unsure{\overline W_n(G)|W_n(G)}}\ \tau:\overline W_n(G)\longrightarrow G_{n-1}}
\end{aligned}
\]

\[
\begin{aligned}
&\deleted{\varphi_n(x_1,\ldots,x_n)=\frac1n\left(1+\frac{p_1}{n-1}
+\frac{p_2}{\frac{(n-1)(n-2)}2}\,+\cdots+
\frac{p_k}{\binom{n-1}{k}}+\cdots+p_{n-1}\right)}\\
&\deleted{\text{où}\quad p_k=\deleted{\text{somme des}}\ \sum x_1 x_2\ldots x_k.}
\end{aligned}
\]

\[
\begin{aligned}
&\deleted{U_j f=\frac{1}{2\pi i}\int_{B_j} f(z_1,\ldots,z_{j-1},t,z_{j+1},\ldots,z_n)\,\frac{dt}{t-z_j}}\\
&\deleted{T_j f=\frac{1}{2\pi i}\iint_{K_j} f(z_1,\ldots,z_{j-1},t,z_{j+1},\ldots,z_n)\,\frac{dt\,d\bar t}{t-z_j}}
\end{aligned}
\]

\deleted{\deleted{\unsure{On|Et}} Pour \deleted{toute}toute suite strictement croissante $j_1<j_2<\ldots<j_q$, on \deleted{\unsure{note}} \inserted{introduit l'opérateur}
$\varphi_n(U\,;j_1,\ldots,j_q)\deleted{=\unsure{\psi_n(}}\inserted{=}$ ce que devient $\psi_n(U_1,\ldots,U_n)$ quand on remplace $U_{j_1}$, \ldots et $U_{j_q}$ par \uline{zéro}.}

\deleted{L'opérateur d'homotopie, pour $\omega=a\,d\bar z_{j_1}\wedge\cdots\wedge d\bar z_{j_q}$ \quad $(q\geq1)$ est}
\margin{\deleted{($n$ est le nombre des variables)}}
\[
\deleted{\underline{\left|\vphantom{\sum^q}\right. P\omega=\varphi_n(U\,;j_1,\ldots,j_q)\,\deleted{\unsure{\sum|\ast}}\sum_{\unsure{\nu|v|\gamma}=1}^{q}(-1)^{\nu-1}
(T_{j_\nu}a)\,d\bar z_{j_1}\wedge\cdots\wedge\widehat{d\bar z_{j_\nu}}\wedge\cdots\wedge d\bar z_{j_q}}}
\]

\[
\deleted{P d'' f\ \deleted{d''P f}=f-U_1 U_2\ldots U_n f}
\]
% archive sheet number: 187
\scanpage{189}
\uline{Données de ramification}

\uline{Catégories multigaloisiennes}
% archive sheet number: 188
\scanpage{191}
\mspage{1}
$S$ schéma, \inserted{$Et(S)$ le site étale de $S$. $\mathcal U$ un univers.}
Une ``donnée de ramification'' sur $S$ est un couple $(\mathscr R,r)$ formé d'un \uline{champ en \inserted{$\mathcal U$-}catégories multigaloisiennes} $\mathscr R$ sur $Et(S)$,
\margin{on sous-entend que les \unsure{fonct$^{\rm s}$|foncteurs} chgt de base sont \uuline{exacts} !}
et d'un morphisme \deleted{de champs} \inserted{cartésien} (``réalisation géométrique'')
\[
r:\mathscr R\longrightarrow \mathrm{Fl}_S
\]
où $\mathrm{Fl}_S$ est la \deleted{\unsure{champ}} catégorie fibrée sur $Et(S)$ dont la fibre en l'objet $S'$ est $(Sch)_{/S'}$ ;
% "catégorie fibrée" is struck through and dotted underneath (stet).
\deleted{\unsure{on donne}} un objet $X$ de $\mathscr R_{S'}$ ( \deleted{\unsure{Et|objet}} $S'\in\mathrm{Ob}\,Et(S))$ sera dit \uline{revêtement de type $\mathscr R$} de $S'$, \deleted{\unsure{et}} $r(X)$ est sa \uline{réalisation} géométrique.
\deleted{On fait l'hypothèse} Ces données sont soumises aux axiomes suivants :

R1 Le foncteur $r$ commute aux sommes finies et au passage au quotient
\inserted{il suffit passage au quotient par groupe fini \unsure{ordinaire} (et \unsure{même|dans les} cas galoisien \unsure{ss pt fixe !})}
par une relation d'équivalence
[il sera $=$ exact à droite, cf \uline{plus bas} \unsure{\ldots}, \uline{mais en général pas exact à gauche} !].
\margin{On veut que dans $\mathrm{Fl}_S$ on ait quotient \uuline{admissible}}

R2 \unsure{Pour tt|\deleted{Donné} Soient} $S'\in\mathrm{Ob}\,Et(S)$ et $X\in\mathrm{Ob}\,\mathscr R_{S'}$, l'application $X'\mapsto r(X')$ est une bijection de l'ens des \deleted{sous-objets}\inserted{sommandes} directs \inserted{(= ss-objets)} $X'$ de $X$ avec l'ens des sommandes directs (i.e. parties ouvertes et fermées) de $r(X)$.

R3 Soit $S'$, $X$ comme dessus, et soit \deleted{\unsure{$(X_i)$}}\inserted{$(X_i)$} une famille (pas néc. finie) de \deleted{\unsure{flèches de but $X$}} sous-objets de $X$, telle que la famille des $r(X_i)$ recouvre $r(X)$. Alors la \deleted{famille $X_i$ admet $X$ comme somme directe effective,}\deleted{\unsure{soit la}} famille $X_i\longrightarrow X$ est une famille de descente effective (= est épimorphique effective \ldots).
% "sous-objets de X" is struck and dotted underneath.

R4 Soit $X\to Y$ dans $\mathscr R_{S'}$ ($S'\in\mathrm{Ob}\,Et(S)$), d'où $Y=\coprod_{n\in\mathbb N}Y_n$, avec $X\times_Y Y_n\to Y_n$ de rang $n$. On veut que les $r(Y_n)$ \uline{recouvrent} $r(Y)$.
\margin{Dans R4, il suffirait de supposer $Y=e$}

\uline{Prop}. \deleted{\unsure{Cor$^{\rm e}$|Écri}} Disons qu'une famille de flèches \inserted{$X_i\longrightarrow X$} de $\mathscr R_{S'}$ est couvrante si la famille $r(X_i)\longrightarrow r(X)$ est surjective, i.e. (grâce : R1 c'est kif kif) si les images $Y_i\subset X$ sont telles que les $r(Y_i)$ recouvrent $r(X)$. Alors ce sont les familles couvrantes d'une \uline{topologie} sur $\mathscr R_{S'}$, \inserted{moins fine que la canonique} \inserted{et pour laquelle tout objet de $\mathscr R_{S'}$ est loc. constant fini.} le foncteur chgt de base $\mathscr R_{S'}\longrightarrow\mathscr R_{S''}$ est continu. \unsure{(donc|Comme|[donc}
\margin{n'utilise que R1 et R3, R4}
% archive sheet number: 190
\scanpage{192}
\mspage{2.}
définit un morphisme de sites en sens inverse).

Axiomes d'une topologie : a) Si une famille est \unsure{majorisée|majorée} par une famille couvrante, elle est couvrante : trivial

b) Un \unsure{isom|somme} est couvrant : trivial \quad c) stabilité par chgt. de base : on est ramené : des $X_i-X$ qui sont des monomorphismes, grâce : b).
\diagram{\begin{tikzcd}[ampersand replacement=\&]
X_i\arrow[d] \& X'_i\arrow[l]\arrow[d]\\
X \& X'\arrow[l]
\end{tikzcd}}
mais alors (mettre en lemme) $r$ commute aux produits fibrés $X_i\times_X X'$, et on gagne

d) Notion locale : supposons qu'on ait des $X'_\alpha-X$ couvrant $X$, tels que $\forall\alpha$, les $X_i\times_\alpha X'_\alpha=X'_{i\alpha}$ couvrent $X_\alpha$. Alors les $X_i$ couvrent $X$. Ici encore, on est ramené au cas où les $X_i\to X$ sont des monomorphismes, et on gagne.
\diagram{\begin{tikzcd}[ampersand replacement=\&]
X_i\arrow[d,two heads] \& X'_{i\alpha}\arrow[d,two heads]\\
Y_i\arrow[d,hook] \& Y'_{i\alpha}\arrow[d,hook]\\
X \& X'_\alpha\arrow[l]
\end{tikzcd}}
(Question : y a-t-il une référence à SGA 4 : image inverse d'une topologie ? Mais attention au fait que $r$ n'est pas exact à g.!). On n'a utilisé pour l'instant que R1. R3 \deleted{implique}\inserted{signifie} que la topologie est \deleted{\unsure{plus}}\unsure{moins} fine que la canonique, R4 que \inserted{les objets de $\mathscr R_{S'}$ sont loc. cts finis}. L'assertion sur les changements de base est triviale.

\uline{Corollaire} Prenant les topos \deleted{\unsure{assoc|associés aux}} $\widetilde{\mathscr R}_{S'}$ associés aux $\mathscr R_{S'}$, on trouve sur $Et(S)$ un \uline{champ en topos}. $\mathscr R_{S'}$ ``sont'' la sous-catégorie des objets loc. constants finis de $\widetilde{\mathscr R}_{S'}$

\deleted{Soit $X\to Y$ un morphisme de $\mathscr R_{S'}$. Alors $Y=\coprod_{n\in\mathbb N}Y_n$, \inserted{(somme \uuline{effective}} avec $X\times_Y Y_n\to Y_n$ \uline{de rang} $n$. \deleted{\unsure{Quelle}} \inserted{La famille} \deleted{\unsure{somme}} $Y_n\to Y$ est elle \uuline{couvrante} ?? \inserted{En d'autres termes,} Les $r(Y_n)\to r(Y)$ couvrent-ils ??}
% archive sheet number: 191
\scanpage{193}
\mspage{3}
\margin{Utilise R1 et R2}
\uline{Prop 2} Soit $X\in\mathrm{Ob}\,\mathscr R_{S'}$. Alors $X=\varnothing\Longleftrightarrow r(X)=\varnothing$, $X$ connexe $\Longleftrightarrow r(X)$ connexe.

Soit $C$ une \inserted{$\mathcal U$-}\deleted{\unsure{cat.}}catégorie multigaloisienne, $P$ un objet.
% "catégorie" is crossed but has a dotted reinstatement.
Un objet $x$ de $C$ est dit \uline{subordonné} à $P$ si $X_P=\deleted{\unsure{P\times X}}X\times P\longrightarrow P$ est \deleted{\unsure{loc. constant pour la top}}\inserted{tel que, sur le ss-objet $P_n$ de $P$ où $X_P$ est de rang $n$, \deleted{$X$} est \uline{constant}.}
\diagram{\begin{tikzcd}[ampersand replacement=\&]
X\arrow[d,dash] \& X\times P\arrow[l,dash]\arrow[d,dash]\\
e \& P\arrow[l,dash]
\end{tikzcd}}
Pour $P$ fixé, les objets $X$ de $C$ subordonnés à $P$ forment une ss-catégorie pleine stable par $\mathop{\mathrm{lim}}\limits_{\substack{\longleftarrow\\\longrightarrow}}$ finies et par sommande direct [car les objets \inserted{sur $P$} qui sont ``constants par morceaux'' forment une \inserted{ss-}catégorie pleine de $C_{/P}$ stable par $\mathop{\mathrm{lim}}\limits_{\substack{\longrightarrow\\\longleftarrow}}$ finies et par sommandes directes \ldots],
\margin{dégager en lemme de catégorie multigaloisienne}
% slanted
donc c'est une ss-catégorie multigaloisienne pleine de $C$, soit $C^{(P)}$. Si $P$ est galoisien de groupe $G$, cette catégorie est équivalente à celle des objets \inserted{dans $P$} cts par morceaux, \deleted{\unsure{munis}} avec op. de $G$ dessus compatible avec son op. sur $P$.

Si maintenant $C=\mathscr R_S$, on peut regarder la ss-catégorie fibrée pleine de $\mathscr R$ formée des $X'\in\mathrm{Ob}\,\mathscr R_{S'}$ qui sont \uuline{localement} \uline{subordonnés} à $P$. On voit encore que c'est stable par $\mathop{\mathrm{lim}}\limits_{\substack{\longrightarrow\\\longleftarrow}}$ et sommandes directes, donc c'est un ss-\deleted{catégorie}\inserted{champ en} \deleted{fibrée}\inserted{catégories multigaloisiennes} plein $\mathscr R^P$ de $\mathscr R$.

Supposons maintenant $P$ galoisien de groupe $G$, alors les objets de $\mathscr R_S$ \uline{subordonnés} (pas seul$^{\rm t}$ loc$^{\rm t}$ !) à $P$ correspondent aux objets sur $P$ \uline{constants} par morceaux, avec op. de $G$ dessus. Or, grâce à R1, R\unsure{4|3|\deleted{2}4}, un objet \inserted{$X$} sur $P$ constant par morceaux définit un objet \inserted{$r(X)$} sur $r(P)$ constant par morceaux, et
\margin{schémas}
% archive sheet number: 192
\scanpage{194}
% Large diagonal strokes cancel both paragraphs. The last line is struck separately.
\deleted{\uline{Proposition} Supposons R1 R3 R4. Alors R2 équivaut à la condition que $r$ est \uline{fidèle}.}

\deleted{Supposons $r$ fidèle, prouvons R2. Soient $X'$, $X''$ deux ss-objets de $X\in\mathrm{Ob}\,\mathscr R_{S'}$ tels que $r(X')=r(X'')$ comme sous-objet de $r(X)$ ; \deleted{donc $r(X')\unsure{\cap}$} or $r(X'\cap X'')=r(X')\cap r(X'')$, (cf lemme) donc l'inclusion $\deleted{\unsure{r(|\text{et}}}X'\cap X''\subset X'$ et $X'\cap X''\subset X''$ \unsure{induit|induit un|deviennent des} isom sur $r$. On est ramené au cas où on a \inserted{déjà} une inclusion $X'\subset X''$, i.e. $X''=X'\amalg X_1$. Alors $r(X'')=r(X')\amalg r(X_1)$, \deleted{et l'hyp: $r(X')=r(X'')$ équivaut à $r(X_1)=\varnothing$.}}
% archive sheet number: 193
\scanpage{195}
\mspage{4.}
\deleted{\unsure{en vertu de}} utilisant de plus R2, on voit que le foncteur $X\unsure{-|\mapsto}r(X)$ des \deleted{\unsure{rev.}} objets sur $P$ cts par morceaux dans les schémas sur $r(P)$ cts par morceaux est pl. fidèle. Utilisant de plus R2 et R3, on trouve qu'il est \unsure{ess surjectif|un surjectif}. C'est donc une équivalence de catégories,
\margin{lemme à dégager ; n'utilise pas l'\unsure{hyp.|hyp} galoisienne sur $P$}
% slanted
\deleted{\unsure{ici|i.e.}} Notons que comme $G$ opère sur \unsure{$P$|$B$|$R$}, il opère sur $r(P)$, \deleted{donc} et on trouve que l'équivalence est compatible avec les opérations de $G$. On trouve ainsi :

\uline{Prop} 3 \inserted{Soit $P$ objet galoisien de groupe $G$ \unsure{de|dans} $\mathscr R_S$.}
Le foncteur $X\unsure{\leadsto|\mapsto}r(X\times P)/r(P)$ induit une \uline{équivalence} de la catégorie des objets de $\mathscr R_S$ subordonnés à $P$, avec la catégorie des schémas \inserted{finis} constants par morceaux sur $r(P)$, sur lesquels $G$ opère de façon compatible avec ses opérations sur $r(P)$.

Utilisant la théorie de la descente, il vient

\uline{Th.} 1 Soit $P$ objet galoisien de groupe $G$ en $\mathscr R_S$. Le foncteur $X\longmapsto r(X\times P)/r(P)$ induit une équivalence de la catégorie des objets de $\mathscr R_S$ loc. subordonnés à $P$, avec la catégorie des \uline{revêtements étales} \inserted{\deleted{\unsure{$Z$}}$Z'$} de $Z=r(P)$, \uline{à action de} $G$ compatibles avec son action sur $r(P)$, et satisfaisant la condition suivante (\unsure{automatiqu$^{\rm t}$|automatique|automatiquement} satisfaite si $r(P)=Z$ est fini sur $S$) : les $S'\unsure{-|\to}S$ étales tels que $Z'_{S'}$ soit \deleted{\unsure{étale sur|étales sur}} constant par morceaux sur $Z_{S'}$ \uline{recouvrent} $S$. / Localisant sur $S$, on trouve de même une description \deleted{de}\inserted{du \uline{champ}} $\mathscr R^P$ en termes du schéma $Z$ à groupe d'opérateurs $G$ sur $S$.
% archive sheet number: 194
\scanpage{196}
\mspage{5.}
Notons d'ailleurs

\uline{Prop. 4} Pour tt schéma \deleted{à opérateurs} $Z$ sur \uline{$S$}, à groupe fini $G$ d'opérateurs, \deleted{\unsure{opé}} agissant de façon admissible, le champ \deleted{défini}\inserted{décrit} \deleted{par la pr}\inserted{comme} cidessus est un champ en \inserted{$\mathcal U$-}catégories multigaloisiennes, et le foncteur ``passage au quotient par $G$'' (qui est défini, car l'action de $G$ reste admissible) \deleted{est} \unsure{font|fait} une donnée de ramification sur $S$.

Tout $\mathscr R_{S'}$ est multigaloisien : tout d'abord, il en est ainsi de la catégorie des \deleted{objets} rev. à opérateurs $G$ \deleted{\unsure{du|des}} $Z_{S'}$, or on en prend ici une sous-catégorie stable par $\mathop{\mathrm{lim}}\limits_{\substack{\longrightarrow\\\longleftarrow}}$ finies et par sommandes directs ; O.K. Foncteurs de chgt de base exact : O.K.

Le foncteur $r$, étant un passage au quotient, commute aux $\varinjlim$ finies (qui ici se calculent au sens de $(Sch)$ !), d'où R1. De $=$ R2 est immédiat. R3 aussi. R4 aussi. O.K.

\deleted{Ainsi} Supposons pour simplifier $S$ \unsure{qu cpct|qu. cpct}, et que $\mathscr R$ satisfasse une petite condition supplémentaire (tjrs vérifiée en pratique) : \deleted{\unsure{tout obj}}

Pour tt objet de $\mathscr R_S$, son rang sur l'objet final de $\mathscr R_S$ \deleted{est}\inserted{reste} borné, i.e $X$ est subordonné à un \deleted{objet} objet à opérateurs galoisien $(P,G)$ de $\mathscr R_S$.

Alors $\mathscr R_S$ est limite inductive \uline{filtrante} de ses sous-catégories de la forme $\mathscr R_S^P$. On a donc une assez bonne description de $\mathscr R_S$, grâce au théorème, en termes d'un système \uline{projectif} d'objets à opérateurs $(Z_i;G_i)$.
% archive sheet number: 195
\scanpage{197}
\mspage{6.}
correspondant à des $(P_i,G_i)$.

Dans le cas où l'objet final de $\mathscr R_S$ est connexe, on peut prendre les $P_i$ connexes, donc les morphismes de transition $G_j-G_i$ sont surjectifs, et on trouve $Z_i\simeq Z_j/H_{\unsure{ij|j}}$ \ldots. D'ailleurs, les $Z_i/G_i$ seront connexes \ldots. D'autre part, relativement à un tel système projectif, on peut répéter un énoncé comme prop. 4, soit

\uline{Prop} 5 \ldots..

Mais \deleted{\unsure{par les}} \uline{partant} de $\mathscr R$, et construisant un système projectif $(Z_i,G_i)$, \deleted{d'où} grâce à celui-ci un nouveau \deleted{champ} $(\mathscr R',r')$, et $\mathscr R'\hookrightarrow\mathscr R^{\deleted{\unsure{\prime|r|\bullet}}}$, il n'est pas vrai en général, \inserted{a priori,} que $\mathscr R'\hookrightarrow\mathscr R$ soit une équivalence : \uline{Le champ} $\mathscr R$ \uline{n'est pas nécessairement} ``\uline{engendré}'' par ses sections globales. Donc pour définir un $\mathscr R$ en \uline{général}, il faut se localiser et recoller \ldots. On va donc décrire un procédé de recollement qui nous permettra de construire les données de ramification dont nous aurons besoin.

Notions auxiliaires : Si $(\mathscr R,r)$, $(\mathscr R',r')$ sont deux données de ram. sur $S$, un hom. de l'une sur l'autre est un couple $(\varphi,\lambda)$, avec $\varphi:\mathscr R\longrightarrow\mathscr R'$ \deleted{\unsure{un}} foncteur cartésien, exact \inserted{sur les fibres,} et $\lambda$ un \uline{isomorphisme} de foncteurs cartésiens $r'\varphi\simeq\deleted{\unsure{\varphi}}r$. \deleted{Si $\mathscr R$ est défini} De cette façon,
% archive sheet number: 196
\scanpage{198}
\mspage{7}
les données de ramification sur $S$ forment une 2-catégorie \inserted{$\unsure{\mathrm{Don\,ram}|\mathrm{Don\,rem}|\mathrm{Domram}}(S)$.} Pour $S'$ variable \inserted{dans $Et(S)$}, les 2-catégories ainsi attachées $\mathrm{Don\,ram}(S')$ forment une 2-catégorie fibrée \ldots.

Si $(\mathscr R,r)$ est défini par $Z,G$, de sorte qu'on a un objet galoisien canonique $P$ dans $\mathscr R_S$, de groupe $G$, un \deleted{foncteur} hom $(\mathscr R,r)-(\mathscr R',r')$ définit un objet $\varphi(P)$ de $\mathscr R'$, qui est galoisien \deleted{de} pour $G$, et un isomorphisme \uline{compatible avec} $G$
\[
r(\varphi(P))\overset{\alpha}{\simeq}\deleted{\unsure{r(P)}} Z
\]

\uline{Prop} 6 On trouve ainsi une équivalence entre $\underline{\mathrm{Hom}}((\mathscr R,r),(\mathscr R',r'))$ et la catégorie des couples $(P',\alpha)$, où $P'$ est un objet galoisien de $\mathscr R^{\unsure{\prime|\,}}_S$, de groupe $G$, et $\alpha$ un $G$-isomorphisme.
\[
r(P')\simeq Z
\]

[En effet, $P',\alpha$ étant donnés, on a une \uline{équivalence} $\mathscr R\cong{\mathscr R'}^{P'}$. On trouve en passant \unsure{ici|\S} :

\uline{Cor} \inserted{Pour} Un hom $(\varphi,\lambda)$ de données de ramification, \deleted{\unsure{on supp}} $\varphi$ est néc. \uline{pleinement} \uline{fidèle}. ]

En particulier

\uline{Cor} \deleted{\unsure{Cor|En}} Si $(\mathscr R',r')$ est défini par $Z',G'$, alors $\underline{\mathrm{Hom}}((\mathscr R,r),(\mathscr R',r'))$ est équivalent à la \uline{catégorie} des fibrés principaux $P'$ de $Z'$, de groupe $G$, sur lesquels $G'$ opère \deleted{\unsure{sur $P'$}} de façon compatible avec ses opérations sur $Z'$, \uline{munis} d'un isom de $G$-schémas sur $S$
\margin{$P'$ \uline{constant par morceaux} loc. (ét) sur $S$,}
% margin note written sideways
\[
P'/G'\simeq Z.
\]
% archive sheet number: 197
\scanpage{199}
\mspage{8.}
On explicite facilmt la composition des hom des données de ramification associées à $Z,G$, $Z',G'$, $Z'',G''$ \ldots.. Ceci dit :

Soient des $S'_i$ couvrant $S$ ; sur chaque $S_i$ supposons donné un $(Z_i,G_i)$ schéma à groupe d'opérateurs $G_i$ admissible ; pour tt couple $(i,j)$, on suppose donné \deleted{\unsure{un}} sur $S_i\times_S S_j$ une ``\uline{subordination}'' de $(Z_{ij},G_i)$ à $(Z_{ji},G_j)$ $(Z_{ij}=Z_i\times_{S_i}(S_i\times_S S_j))$, \deleted{de telle façon qu'on ait} sur les $S_{ijk}$ des isomorphismes de transitivité entre les subordinations. On suppose enfin une condition de compatibilité sur les $S_{ijkl}$. Alors on peut recoller \deleted{les} $(\mathscr R_i,r_i)$ définis par les $(Z_i,G_i)$ \ldots..

Donnons un cas particulier moins enclin à donner le mal de mer :

Soit une catégorie fibrée \inserted{(de ``modèles'')} $C$ sur $Et(S)$ (pas néc. un champ) et un foncteur cartésien $s:C\unsure{\longrightarrow|-}\mathrm{Fl}_S$. On suppose $C$ fibrée en groupoïdes, que deux objets d'un $C_{S'}$ sont localmt isomorphes, que $C_S$ est localement non vide -- d'où finalement, en prenant son champ associé, une \uline{gerbe}. On suppose que le lien de cette gerbe est défini \deleted{\unsure{par un objet}} loc. constant tordu, i.e. pour tt $x\in C_{S'}$, le faisceau associé
% archive sheet number: 198
\scanpage{200}
\mspage{9.}
au préfaisceau $S''\longrightarrow\mathrm{Aut}\,x_{S''}$ est représentable par un \deleted{objet groupe de la catégorie des} rev. étale de $S'$. On a une ``représentation'' de cette gerbe dans le préchamp $\mathrm{Fl}_S$. Á ces \unsure{belles|telles} données, on associe une donnée de ramification sur $S$ ainsi

Les objets de $\mathscr R_S$ sont les hom. cartésiens \uline{sur $\mathrm{Fl}_S$} $\mathscr G\longrightarrow\mathrm{Rev}\,\mathrm{Fl}_S$ qui relèvent $\mathscr G\unsure{-|\to}\mathrm{Fl}_S$, et itou quand $r$ est remplacé par $r'$.
\diagram{\begin{tikzcd}[ampersand replacement=\&]
{}\&{}\&\mathrm{Rev}\,\mathrm{Fl}_S\arrow[d]\\
\mathscr G\arrow[rr]\arrow[dr,dash]\&{}\&\mathrm{Fl}_S\arrow[dl]\\
{}\&Et(S)\&{}
\end{tikzcd}}
Le foncteur $r_S:\mathscr R_S\longrightarrow (Sch)_{/S}$ est le foncteur $\varinjlim$. On vérifie que tt marche.

\uline{Ex 1} Donnons \deleted{\unsure{la}} famille $\underline a=(a_i)_{i\in\unsure{I|J}}$ de sections de $\mathcal O_S$, et famille $\underline n=(\underline n_i)_{i\in I}$ d'entiers $\geq1$, \inserted{premiers aux car. résiduelles} \deleted{\unsure{les points}} et tous égaux à 1 sauf un nb fini. \deleted{\unsure{tout sauf}} Soit $C_{S'}=\deleted{\unsure{F\!l}}$ groupoïde défini par le groupe $\mu_{\underline n}(S')=\prod_{i\in I}\mu_{n_i}(S)$, \deleted{Supposons}\inserted{(unique objet $e_{S'}$)} et $s(e_{S'})=Z_{\underline n}^{\underline a}$, avec l'opération qu'on devine de ce groupe. Les axiomes sont vérifiés. [En fait, cela correspond au cas d'une \uline{gerbe triviale} définie par $Z_{\underline n}^{\underline a}$ sur lequel $\mu_{\underline n}^{\underline a}$ opère ].

\uline{Ex 1 bis} \inserted{Donnée \deleted{de} de ram. Kummer définie par $\underline a,\underline n$.} Comme avant, mais on passe à la limite sur $\underline n$ variable, $\underline a$ fixé ; donnée de ramification kummérienne définie par $\underline a$.

\uline{Ex} 2 Donnons famille $(D_i)_{i\in I}$ de diviseurs, et famille $(n_i)_{i\in I}$ d'entiers satisfaisant mêmes conditions que dans ex. \uline{1}. On prend
% archive sheet number: 199
\scanpage{201}
\mspage{10.}
$C_{S'}$ défini ainsi :
\[
\begin{aligned}
\mathrm{Ob}(C_{S'})&=\{\underline a=(a_i)_{i\in\unsure{I|J}}\mid \forall i,\ a_i\ \text{une équation de }D_i\}\\
\mathrm{Hom}(\underline a,\underline b)&=\{\unsure{\zeta|\xi}=(\zeta_i)_{i\in\unsure{J|I}},\mid\forall i,\ b_i=a_i\zeta_i^{n_i}\}
\end{aligned}
\]

\deleted{Image inverse et} Composition des hom évidente. Images inverses aussi. On trouve comme champ associé une gerbe abélienne de groupe $\mu_{\underline n}$. (L'existence d'une section \uline{dépend} de la nullité d'un élément de $H^2(X,\mu_{\underline n})$, qui n'est autre, comme on devine, que la famille des classes de cohomologie des $D_i$ dans les $H^{\unsure{2|1}}(X,\mu_{n_i})$ !) : \uline{donnée de ramification kummérienne} définie par $(D_i)_{i\in I}=D$ et $\underline n$.

\uline{Ex 2 bis} On donne $D=(D_i)_{i\in I}$, on passe à la limite sur $\underline n$ : donnée de ramification kummérienne définie par $D$.

Ces exemples boitent un tout petit peu, on voudrait en fait que $n_i$ soit premier aux car résiduelles de $D_i$ seulement, pas néc. celles de tt $S$. Dans le cas de l'exemple 1, on prendra les rev. étales de $Z_{\underline n}^{\underline a}$, sur lesquels $\mu_n$ opère de façon compatible avec ses opérations sur $Z_{\underline n}^{\underline a}$. Dans le cas de l'exemple 2
% archive sheet number: 200
\scanpage{202}
\mspage{11.}
Soit $(\mathscr R,r)$ une donnée de ramification sur $S$. Une famille d'objets $X_i/S_i$ de $\mathscr R$ ($S_i\in\mathrm{Ob}\,Et(S)$) est dite \uline{génératrice} si pour tt objet $X$ de $\mathscr R$ sur un $S'$ de $Et(S)$, les $S''$ étales sur $S$ tels qu'il existe $i$ et $S''\unsure{-|\to}S_i$ avec $X_{S''}$ subordonné à $X_{i\,S''}$ \inserted{et $X_{i\,S''}\deleted{\unsure{\to}}\to e_{S''}$ \deleted{\unsure{étant}} couvrant}, recouvrent $S'$.
\diagram{\begin{tikzcd}[ampersand replacement=\&]
S\arrow[dr,dash]\&{}\&S_i\\
{}\&S'\arrow[dr,dash]\&{}\\
{}\&{}\&S''
\end{tikzcd}}
Si \inserted{de plus} les objets $X_i$ sont galoisiens dans $\mathscr R_{S_i}$, \deleted{\unsure{sur les}} de groupe $G_i$, d'où $(Z_i,G_i)$ objet à opérateurs $r(\deleted{\unsure{P|R}}X_i)$ sur $S_i$, les $Z_i,G_i$ doivent permettre de \uline{reconstruire} $(\mathscr R,r)$, une fois précisé des données supplémentaires permettant de recoller ; mais la formulation de telles données est pénible et finalement peu utile.

Il est plus commode de supposer donnée déjà une \uline{catégorie fibrée} $C$ sur \inserted{$Et$}$(S)$, et un foncteur \inserted{cartésien} $\varphi:C\longrightarrow\mathscr R$. On dit que $\varphi$ est générateur si les $\varphi(M)$ ($M\in\mathrm{Ob}\,C$) forment une famille génératrice, i.e si pour tt $X\in\mathrm{Ob}\,\mathscr R$ sur un $S'\in\mathrm{Ob}\,Et(S)$, les $S''$ étales sur $S'$ tels que $X_{S'}$ soit subordonné à un objet \inserted{$M$} de $C_{S''}$ \inserted{avec $\varphi(M)\to e_{S^{\unsure{\prime|\prime\prime}}}$ couvrant,} forment un recouvrement de $S'$. Posant $s=r\varphi$
\[
s:C\unsure{-|\longrightarrow}\mathrm{Fl}_S
\]
je dis qu'on peut alors, à équivalence près,
% archive sheet number: 201
